\begin{slide}
    \begin{block}{Problem}
		\begin{itemize}
			\item Adam placed his electronics randomly into his room.
			\item Not every device can be plugged into every socket.
			\item How many devices can Adam power using 1 powerbar which triples 1 socket?
		\end{itemize}
    \end{block}
    \begin{block}{Solution}
		\begin{itemize}
			\item Naive idea: Compute a maximum bipartite matching (possibly using a maximum flow algorithm) for each of the $n$ possible locations of the powerbar
			\item This is, however, too time-consuming!
		\end{itemize}
    \end{block}
\end{slide}

\begin{slide}
%    \begin{block}{Problem}
%		\begin{itemize}
%			\item Adam placed his electronics randomly into his room.
%			\item Not every device can be plugged into every socket.
%			\item How many devices can Adam power using 1 powerbar which triples 1 socket?
%		\end{itemize}
%    \end{block}
    \begin{block}{Solution}
		\begin{itemize}
			\item Create a bipartite graph with sockets on the left, devices on the right and a connection between a socket and a device if the device can be plugged into the socket.
			\item Compute a maximum bipartite matching without the powerbar
			\item For each matched socket with more than 1 neighbor:
			\begin{itemize}
				\item Triple the socket (i.e. add 2 unoccupied copies of the socket / raise the capacity in the flow network).
				\item Compute the maximum bipartite matching / maximum flow in the resulting graph.
				\item Use the bipartite matching of the original graph as a starting point for the matching of the new graph!
			\end{itemize}
			\item Take the maximum across all computed matchings.
		\end{itemize}
    \end{block}
\end{slide}
